\documentclass[10pt,a4paper]{amsart}
\usepackage[margin=19mm]{geometry}
\usepackage{amsmath,amssymb,mathtools}
\usepackage[scr=rsfs,cal=euler]{mathalfa}
\usepackage{xfrac}
\usepackage[unicode,hidelinks,bookmarks=false]{hyperref}
\newcommand{\ie}{i.e.\ }
\newcommand{\cf}{cf.\ }
\newcommand{\resp}{respectively\ }
\newcommand{\Subsection}{\subsection}
\providecommand{\nobreakdash}{\nobreak\hbox{-}}
\setlength{\emergencystretch}{2em}
\multlinegap0pt
\usepackage{bm}
\usepackage{bbm}
\usepackage{dsfont}
\usepackage{xspace}
\usepackage{cancel}
\usepackage{mathtools}
\usepackage[shortlabels]{enumitem}
\usepackage{amsthm}
\makeatletter
\@ifundefined{theorem}{\newtheorem{theorem}{Theorem}}{}
\@ifundefined{claim}{\newtheorem{claim}[theorem]{Claim}}{}
\@ifundefined{lemma}{\newtheorem{lemma}[theorem]{Lemma}}{}
\@ifundefined{proposition}{\newtheorem{proposition}[theorem]{Proposition}}{}
\makeatother
\newcommand{\var}{\varnothing}
\newcommand{\w}{\omega}
\title[Countably compact inverse semigroups and Nyikos' problem]{Countably compact inverse semigroups and Nyikos' problem}
\author{Serhii Bardyla}
\date{Source: arXiv:2503.13666v2 (CC BY 4.0)}
\begin{document}
\maketitle

\noindent\emph{Source and licence.} Countably compact inverse semigroups and Nyikos' problem. Version arXiv:2503.13666v2 (CC BY 4.0), distributed under the Creative Commons Attribution 4.0 International licence, as recorded in the arXiv licence metadata for this exact version and shown on the abstract page https://arxiv.org/abs/2503.13666v2. Full rights notice: licenses/arxiv-2503.13666-CC-BY-4.0.txt.\par\smallskip

\noindent\textit{Editorial scope.} Counted unit: the Lemma whose source label is \texttt{+1} in the article (source0.tex lines 603--605, source label \texttt{+1}), delivered together with the article's own complete original proof of it (source0.tex lines 607--642, one \texttt{proof} environment of 36 lines). No mathematics is written, repaired or re-derived: statement and proof are the article's own text line for line. Required context retained uncounted: chainfolk (lines 501--503). Every definition, display and label of those lines is retained unchanged. Omitted from this package and not redistributed: 1--500, 504--602, 606--606, 643--1182. Nothing inside a retained excerpt is omitted or altered: every retained body is delivered exactly as the article's lines are written, so no omission receipt of any kind accompanies this package. Citations: the retained excerpts carry no citation command at all, so no bibliography entry of the article is transcribed and no reference is redistributed. Reference adaptation: none was needed and none was made; every reference inside the delivered unit resolves inside it, so no unresolved reference or citation is produced. Printed slips kept literal: no printed word, symbol, hypothesis or number of the counted statement or of its proof was altered. Layout and class: the article's own class is \texttt{article} with packages including amsmath, amsthm, amssymb, enumitem, graphicx, fontenc; the wrapper is the lane's pinned \texttt{amsart} editorial wrapper at 10pt on A4, which restates the theorem-like environments the retained text needs and adds the article's own macro definitions that the retained text needs (\texttt{\textbackslash var}, \texttt{\textbackslash w}), with extra wrapper packages (bm, bbm, dsfont, xspace, cancel, mathtools, enumitem). The delivered numbering is the unit's own. Assets: no retained span contains a figure environment, a graphics-inclusion command, a PDF-page inclusion, a table or a TikZ picture; no image file is copied into this package, the package declares no asset dependency and no image file is redistributed with it. Licence: version arXiv:2503.13666v2 of this article is distributed under the Creative Commons Attribution 4.0 International licence, as recorded in the arXiv licence metadata for this exact version and shown on the abstract page https://arxiv.org/abs/2503.13666v2. A full rights notice is delivered at licenses/arxiv-2503.13666-CC-BY-4.0.txt. Characters: every retained excerpt is pure ASCII, so the delivered engine, XeLaTeX with its default Latin Modern text font, reports no missing character.
\begin{proposition}[Folklore]\label{chainfolk}
The closure of a chain in a semitopological semilattice is a chain.    
\end{proposition}

\begin{lemma}\label{+1}
Let $X$ be a semitopological semilattice and $L\subseteq X$ be a chain isomorphic to $(\alpha,\min)$ for some ordinal $\alpha$. If $\overline{L}$ does not contain maximum, then $L$ is cofinal in $\overline{L}$ and $\overline{L}$ is isomorphic to $(\alpha,\min)$.    
\end{lemma}  

\begin{proof}
Let $L=\{l_\xi:\xi\in\alpha\}$, where $l_\xi\leq l_{\eta}$ if and only if $\xi\leq \eta$. By Proposition~\ref{chainfolk}, $\overline{L}$ is a chain.  Fix any $x\in\overline{L}$. If $L\subseteq {\downarrow}x$, then $\overline{L}\subseteq {\downarrow}x$, as ${\downarrow}x$ is closed in $X$. But then $x$ is the maximum of $\overline{L}$, a contradiction. So, for each $x\in\overline{L}$, the set $L\setminus {\downarrow}x$ is nonempty. It follows that the ordinal $\alpha$ is limit and $L$ is cofinal in $\overline{L}$.
For every $x\in \overline{L}\setminus L$ let $\delta(x)=\min\{\xi\in\alpha: x\leq l_\xi\}$. Since $L$ is cofinal in $\overline{L}$, the ordinals $\delta(x)$, $x\in \overline{L}\setminus L$, are well defined. 

Seeking a contradiction, assume that there exists an infinite decreasing chain $\{x_n:n\in\w\}\subseteq \overline{L}$. The chain $L$, being isomorphic to an ordinal, contains only finite decreasing chains. Hence we lose no generality assuming that $\{x_n:n\in\w\}\subseteq \overline{L}\setminus L$. Observe that $\{\delta(x_n):n\in\w\}$ is a nonincreasing chain in $\alpha$. Thus there exists $k\in\w$ such that $\delta(x_n)=\delta(x_m)$ for each $n,m\geq k$. Then the set $W=X\setminus ({\uparrow}x_k\cup {\downarrow}x_{k+2})$ is a neighborhood of $x_{k+1}$ which is disjoint with $L$. The obtained contradiction implies that $\overline{L}$ contains only finite decreasing chains. Hence $\overline{L}$ is isomorphic to $(\theta,\min)$ for some ordinal $\theta$. Fix the order isomorphism $\psi:\overline{L}\rightarrow \theta$. It is easy to see that $\psi$ can be defined  recursively as follows: $\psi(l_0)=0$ and for each $y\in\overline{L}$,
$$\psi(y)=
\begin{cases}
 \psi(x)+1, & \hbox{if }\exists x<y \hbox{ such that } \{z\in \overline L: x<z<y\}=\varnothing;\\
 \sup\{\psi(x):x< y\}, & \hbox{otherwise}.
\end{cases}
$$

Since $\psi(l_\xi)\geq \xi$ for all $\xi\in\alpha$, we get $\alpha\leq \theta$. 


\begin{claim}\label{claimineq}
$\psi(l_{\xi})\leq \xi+1$ for every $\xi<\alpha$.    
\end{claim}

\begin{proof}
Assume that $\psi(l_\delta)<\delta+1$ for all $\delta<\xi$. There are two cases to consider:
\begin{enumerate}[label=(\roman*)]
    \item $\xi=\mu+1$ for some ordinal $\mu<\alpha$;
    \item $\xi$ is a limit ordinal.
\end{enumerate}

In case (i) we have $\psi(l_\mu)\leq \mu+1=\xi$. Observe that the set $A=\overline{L}\setminus({\uparrow}l_\xi\cup{\downarrow}l_{\mu})$ is open in $\overline{L}$ and disjoint with $L$. Therefore $A=\var$. Hence there exists no $z\in\overline{L}$ such that $l_\mu<z<l_\xi$. The definition of $\psi$ implies that $\psi(l_\xi)=\psi(l_\mu)+1\leq \xi+1$.

Consider case (ii). Let $\pi$ be the supremum of the set $\{l_{\delta}:\delta<\xi\}$ in $\overline{L}$. It is clear that $\pi\leq l_\xi$. Since the set $\{l_\delta:\delta<\xi\}$ is cofinal in $\{x\in\overline{L}:x<\pi\}$ and $\psi$ is an order isomorphism, we get that $$\psi(\pi)=\sup\{\psi(l_\delta):\delta<\xi\}\leq \sup\{\delta+1:\delta<\xi\}=\xi,$$
where the inequality above follows from the inductive assumption. So, if $\pi=l_\xi$, then we are done. Assume that $\pi<l_\xi$. Observe that the set $\overline{L}\setminus({\downarrow}\pi\cup {\uparrow}l_\xi)$ is empty, as it is open in $\overline{L}$ and disjoint with $L$. Hence there exists no $z\in\overline{L}$ such that $\pi<z<l_\xi$. Then 
$\psi(l_\xi)=\psi(\pi)+1\leq \xi+1$. 
\end{proof}

Since the ordinal $\alpha$ is limit, $\alpha=\sup\{\xi+1: \xi<\alpha\}$. Observe that $\psi(L)$ is cofinal in $\theta$, as $\psi$ is an order isomoprhism. Then   
Claim~\ref{claimineq} yields $\theta\leq \sup\{\xi+1: \xi<\alpha\}=\alpha$. Thus $\theta=\alpha$, as required. 
\end{proof}

\end{document}
